% =============================================================================
% General Conclusion & Trilingual Abstract
% Assigned to: Ghaida Farhaoui
% =============================================================================

\chapter*{General Conclusion}
\addcontentsline{toc}{chapter}{General Conclusion}
\markboth{General Conclusion}{General Conclusion}

% TODO: Assigned to Ghaida Farhaoui
% This section synthesizes the project's achievements, academic and technical
% contributions, operational challenges encountered, and perspectives for future work.

\section*{Synthesis of Contributions}
The MindSpark integration project demonstrated the feasibility of pairing low-latency collaborative whiteboards with a decoupled, resilient Generative AI engine. Over three intensive Scrum sprints, our six-student engineering squad successfully designed, implemented, and deployed a cloud-native platform comprising 21 containerized services, fulfilling an ambitious 64-story product backlog.

\section*{Technical and Academic Takeaways}
Building MindSpark provided concrete engineering insights into:
\begin{itemize}[leftmargin=1.5em, topsep=2pt, itemsep=1.5pt, parsep=0pt]
    \item Designing autonomous domain microservices adhering strictly to the \textit{Database-per-Service} pattern using PostgreSQL 16 and Spring Boot.
    \item Managing stateful real-time bi-directional protocols with native WebSockets behind a reactive reverse proxy.
    \item Isolating unpredictable third-party AI APIs through concurrency rate limits and Resilience4j circuit breakers to safeguard core whiteboard responsiveness.
    \item Decoupling event-driven notifications across technology stacks via Apache Kafka in KRaft mode.
    \item Orchestrating complex distributed environments with Docker Compose and monitoring runtime health using Prometheus, Tempo, and Grafana.
\end{itemize}

\section*{Limitations and Perspectives for Future Work}
While the platform provides a complete end-to-end brainstorming workflow, several valuable avenues remain for future development:
\begin{itemize}[leftmargin=1.5em, topsep=2pt, itemsep=1.5pt, parsep=0pt]
    \item Migrating the local canvas synchronization protocol to Conflict-Free Replicated Data Types (CRDTs) to support fully offline-capable, peer-to-peer editing sessions.
    \item Transitioning the single-node Docker Compose setup to a Kubernetes (K8s) cluster with automated horizontal pod autoscaling (HPA).
    \item Incorporating multi-modal vision models to interpret hand-drawn whiteboard sketches directly into structured digital components.
\end{itemize}

\clearpage
% -----------------------------------------------------------------------------
% Trilingual Abstract / Résumé / الملخص
% -----------------------------------------------------------------------------
\markboth{Abstract}{Abstract}

\section*{Abstract}
\noindent In modern remote and hybrid teams, effective brainstorming often stalls due to blank-canvas hesitation, visual clutter, and the tedious manual work needed to turn dozens of scattered notes into a workable action plan. This integration project presents \textbf{MindSpark}, a real-time, AI-assisted collaborative brainstorming platform developed at ISET Bizerte. Engineered as a cloud-native distributed architecture, MindSpark combines an Angular 22 single-page application utilizing reactive Signals and native WebSockets with an ecosystem of eight Spring Boot microservices on Java 21 and a Node.js notification engine. The platform embeds a decoupled AI microservice that executes multi-model generative operations (idea generation, semantic clustering, risk analysis, and action plan synthesis) through Google Gemini and Groq, safeguarded by concurrency limiters and circuit breakers. Resource governance is maintained via Keycloak OAuth2 authentication, a database-per-service PostgreSQL topology, Apache Kafka event streaming, and Polar subscription billing. The entire platform operates as a 21-container Docker Compose orchestration monitored via Prometheus, Grafana Tempo, and Grafana.

\medskip
\noindent\textbf{Keywords:} Collaborative Whiteboard, Real-Time WebSockets, Generative AI, Microservices Architecture, Spring Boot, Angular, Apache Kafka, Keycloak, Docker Compose.

\vspace{0.35cm}

\section*{Résumé}
\noindent Dans les environnements de travail collaboratifs et hybrides actuels, les sessions de remue-méninges (\textit{brainstorming}) se heurtent fréquemment au syndrome de la page blanche, à la surcharge informationnelle et à la lenteur de la synthèse manuelle nécessaire pour transformer des notes éparses en plans d'action structurés. Ce projet d'intégration présente \textbf{MindSpark}, une plateforme collaborative temps réel assistée par intelligence artificielle développée à l'ISET Bizerte. Conçue selon une architecture distribuée orientée microservices, MindSpark associe une application monopage Angular 22 reposant sur les Signals et des WebSockets natifs à un écosystème de huit microservices Spring Boot sous Java 21 et un moteur de notifications Express.js sous Node.js. La plateforme intègre un microservice d'IA découplé exécutant des opérations génératives multi-modèles (génération d'idées, regroupement sémantique, évaluation des risques et synthèse de feuilles de route) via Google Gemini et Groq, protégées par des limiteurs de concurrence et des disjoncteurs (\textit{circuit breakers}). La gouvernance des ressources s'appuie sur Keycloak pour l'authentification OAuth2, une topologie PostgreSQL isolée par service, le streaming d'événements Apache Kafka et la facturation Polar. L'ensemble est orchestré via 21 conteneurs Docker Compose supervisés par Prometheus, Grafana Tempo et Grafana.

\medskip
\noindent\textbf{Mots-clés :} Tableau Blanc Collaboratif, WebSockets Temps Réel, Intelligence Artificielle Générative, Microservices, Spring Boot, Angular, Apache Kafka, Keycloak, Docker Compose.

\vspace{0.35cm}

\section*{Arabic Abstract}
\begin{center}
\fcolorbox{PrimaryNavy!50}{PrimaryNavy!4}{%
  \parbox[c][2.0cm][c]{0.92\textwidth}{%
    \centering\color{PrimaryNavy}
    \textbf{\large [ Arabic Abstract --- To be finalized in Arabic script by Ghaida ]}\\[6pt]
    \small\textit{\color{SlateGray}(Note: Dedicated Arabic typography will be configured for final print compilation)}%
  }%
}
\end{center}
